% Standalone appendix fragment.  Requires amsmath, amssymb, amsthm.
% The only ambient notation is ordinary intuitionistic first-order logic
% with actual equality.  This is a proved alternative induction interface,
% not a claim of a globally smaller complete checker.
\section{A constant-recurrence rule in place of induction}
\label{sec:constancy-with-split}

This section gives a separate, complete arithmetic presentation.  Its
primitive inference rule has neither an induction hypothesis nor an
arbitrary update function: it says that a term which is unchanged by
successor has its initial value everywhere.  The price is an explicit
Boolean equality operation.  That price is part of the presentation and
must be counted in any simplicity comparison.

\subsection{The complete arithmetic base}

Use ordinary intuitionistic first-order logic with actual equality, on
one numerical sort.  Terms have the grammar
\[
 t ::= x\mid 0\mid S(t)\mid \delta(t,t)
          \mid R(n,a;i,u.h).
\]
The last constructor binds $i,u$ in $h$ only.  Its body is finite earlier
term syntax; it is not a variable ranging over numerical functions.
All substitutions are capture avoiding.  Put $1=S0$ and
$s\ne t=(s=t\to\bot)$.  The nonlogical axioms are the universal closures
of the following four displayed schemata:
\begin{align}
 R(0,a;i,u.h)&=a, \tag{R0}\label{eq:const-r0}\\
 R(Sn,a;i,u.h)&=h[n/i,R(n,a;i,u.h)/u],
       \tag{RS}\label{eq:const-rs}\\
 0=1&\longrightarrow\bot, \tag{Sep}\label{eq:const-sep}\\
 \bigl(s=t\land\delta(s,t)=0\bigr)&\ \lor\
 \bigl(s\ne t\land\delta(s,t)=1\bigr).
       \tag{Split}\label{eq:const-split}
\end{align}
The substitution in \textup{(RS)} is simultaneous.  No coverage axiom,
successor-injectivity axiom, induction schema, or equality-reflection
theorem is silently included.

The one additional inference schema is
\[
 \frac{\Gamma\vdash f(Sz)=f(z)}
      {\Gamma\vdash f(t)=f(0)}\quad(\mathrm{Const}),
 \tag{Const}\label{eq:const-rule}
\]
where $f(x)$ is a fixed numerical term template, $z$ is fresh for
$\Gamma,t$ and the passive parameters of $f$, and the displayed premise
contains no new numerical or equational assumption.  The context
$\Gamma$ may contain arbitrary first-order formulas.

The underlying logic can be taken to be standard intuitionistic natural
deduction: assumption; introduction and elimination for
$\land,\lor,\to,\forall,\exists$; falsum elimination; equality reflexivity
and substitution; and the usual eigenvariable conditions for
$\forall$ introduction and $\exists$ elimination.  Thus this presentation
does not use a classical inference rule or a choice principle.

\subsection{Elementary computations require no induction}

Define
\[
 C(v,a,b)=R(v,a;i,u.b),\qquad
 a+b=R(b,a;i,u.Su),\qquad
 P(v)=R(v,0;i,u.i).
\]
Their displayed equations are direct instances of the computation axioms:
\begin{align*}
 C(0,a,b)&=a,& C(Sv,a,b)&=b,\\
 a+0&=a,& a+Sb&=S(a+b),\\
 P0&=0,& P(Sv)&=v.
\end{align*}
In particular $a+1=Sa$ and successor is injective.

Every successor is different from zero.  Indeed, if $Sv=0$, applying
$w\mapsto C(w,0,1)$ gives $1=0$, contradicting \textup{(Sep)}.
This argument uses no case distinction on an arbitrary number.

The split axiom, together with $0\ne1$, gives
\[
 \delta(s,t)=0\lor\delta(s,t)=1,
 \qquad \delta(s,t)=0\ \longleftrightarrow\ s=t.
 \tag{E}\label{eq:const-eqreflection}
\]
For the reverse direction of the equivalence, eliminate the two
disjuncts of \textup{(Split)} under $s=t$; its second disjunct contradicts
that assumption.  For the forward direction, its second disjunct would
give $0=1$.  These are intuitionistic derivations.

\subsection{Boolean failure terms for quantifier-free formulas}

For every quantifier-free formula $A$, recursively define a term $q_A$,
with value zero meaning that $A$ holds:
\begin{align*}
 q_{s=t}&=\delta(s,t), &q_\bot&=1,\\
 q_{A\land B}&=C(q_A,q_B,1),
 &q_{A\lor B}&=C(q_A,0,q_B),\\
 q_{A\to B}&=C(q_A,q_B,0).
\end{align*}
Negation is the abbreviation $A\to\bot$.  Structural induction on the
formula, using the two computations for $C$, proves
\[
 q_A=0\lor q_A=1,
 \qquad q_A=0\ \longleftrightarrow\ A.
 \tag{B}\label{eq:const-boolean}
\]
The induction here is a metatheoretic induction on a finite formula, not
an arithmetic induction inside the proposed theory.  In the implication
case, split on $q_A$: when it is zero, $A$ is available and $A\to B$ is
equivalent to $B$; when it is one, $\neg A$ follows and the implication
is true.  Only the decidability already proved for the quantifier-free
antecedent is used.  No claim is made for formulas with quantifiers.

\subsection{An accumulator turns progress into constancy}

Suppose $A(n)$ is quantifier free and that the proposed theory has
derivations
\[
 \Gamma\vdash A(0),\qquad
 \Gamma,A(z)\vdash A(Sz),
\]
with the usual fresh-variable condition.  Write $q(n)=q_{A(n)}$.
Define the following finite terms:
\begin{align*}
 e(n)&=C\bigl(q(n),0,C(q(Sn),1,0)\bigr),\\
 a(n)&=R(n,0;i,u.u+e(i)),\\
 F(n)&=a(n)+q(n).
\end{align*}
Here $e$ and the addition operation are expanded into their finite term
abbreviations before the outer $R$ is formed.  In particular, the
definition of $a$ does not presuppose a new function symbol or a new
recursion principle.

The computations give
\[
 a(0)=0,\qquad a(Sn)=a(n)+e(n).
\]
The three possible pairs of Boolean values of $q(n),q(Sn)$ give the
following equalities without an induction hypothesis:
\[
\begin{array}{c|c|c|c}
 q(n)&q(Sn)&e(n)&F(Sn)=F(n)\\ \hline
 0&0&0&a(n)=a(n)\\
 1&1&0&Sa(n)=Sa(n)\\
 1&0&1&Sa(n)=Sa(n)
\end{array}
\]
The remaining pair $(0,1)$ is impossible: $q(n)=0$ gives $A(n)$;
the retained step derivation gives $A(Sn)$; and
\eqref{eq:const-boolean} then gives $q(Sn)=0$, contradicting $0\ne1$.
Thus, by finite disjunction elimination under $\Gamma$,
\[
                       \Gamma\vdash F(Sz)=F(z).
\]
Apply \textup{(Const)} to obtain $F(t)=F(0)$.  The proved base $A(0)$
gives $q(0)=0$, so $F(0)=0$.  If $q(t)=1$, then
\[
                 0=F(t)=a(t)+1=Sa(t),
\]
contrary to successor separation.  The Boolean disjunction for $q(t)$
therefore gives $q(t)=0$, and hence $A(t)$.

\begin{theorem}[Recovery of contextual quantifier-free induction]
The four axiom schemata \textup{(R0)}, \textup{(RS)}, \textup{(Sep)},
\textup{(Split)} and the rule \textup{(Const)} derive
\[
 \frac{\Gamma\vdash A(0)\qquad\Gamma,A(z)\vdash A(Sz)}
      {\Gamma\vdash A(t)}
\]
for every quantifier-free $A$, under arbitrary retained first-order
contexts.  In particular they derive contextual induction for every
literal equation.
\end{theorem}

No commutativity or associativity of addition was used in this proof.
The only addition identities used are $a+0=a$ and $a+1=Sa$.
This is why the arithmetic bootstrap does not conceal an induction
argument in an algebraic rearrangement.

\subsection{Coverage, primitive recursion, and both interpretations}

Apply the recovered rule to the quantifier-free formula
\[
                  A(n)\equiv n=0\lor n=S(Pn).
\]
Its zero case is immediate.  At a successor, $P(Sn)=n$ gives the second
disjunct without an induction hypothesis.  Thus numerical coverage is
derived.  Successor injectivity was already obtained from $P(Sn)=n$.

Every usual primitive recursive definition compiles directly into this
counter-aware language: zero, successor, projections, and composition use
the corresponding term operations, and
\[
 f(0,\vec x)=g(\vec x),\qquad
 f(Sn,\vec x)=h(n,f(n,\vec x),\vec x)
\]
is represented by
\[
 R\bigl(n,g(\vec x);i,u.h(i,u,\vec x)\bigr),
\]
after recursively compiling $g,h$.  The computation equations are
literal instances of \textup{(R0)} and \textup{(RS)}.  Together with
successor separation, injectivity, and recovered quantifier-free
induction, this gives an interpretation of full first-order iPRA.

Conversely, each fixed finite $R$ term denotes a primitive recursive
function in iPRA by ordinary primitive recursion on its counter.  Its
two defining equations hold by construction.  Interpret $\delta$ by a
Boolean equality characteristic.  For an explicit construction, define
\[
                         d(a,b)=R(b,a;i,u.Pu).
\]
The subtraction equations and equation induction give
\[
 d(a,0)=a,\quad d(a,Sb)=P(d(a,b)),\quad
 d(Sa,Sb)=d(a,b),\quad d(Sn,n)=1.
\]
For the third identity, induct on $b$ at fixed $a$; the base is
$P(Sa)=a$, and the step applies $P$ to the preceding equality.
For the fourth, induct on $n$ using the third identity.  Elementary
equation inductions on the second argument of addition also give
$0+b=b$ and $Sa+b=S(a+b)$.

Quantifier-free induction on $b$, with $a$ passive, proves
\[
           d(a,b)\ne0\ \longrightarrow\ a=b+d(a,b).
 \tag{Rec}\label{eq:const-reconstruct}
\]
The base uses $0+a=a$.  For the step, put $r=d(a,b)$ and suppose
$Pr=d(a,Sb)\ne0$.  Coverage on $r$ gives $r=S(Pr)$: the zero case
contradicts $Pr\ne0$, and the successor case uses the equation for $P$.
In particular $r\ne0$, so the preceding induction hypothesis gives
\[
              a=b+r=b+S(Pr)=Sb+Pr.
\]
This establishes the new instance of \eqref{eq:const-reconstruct}.

Now define
\[
 Z(a)=C(a,1,0),\qquad U(a)=C(a,0,Z(Pa)),\qquad
 E(a,b)=C\bigl(b,Z(a),U(d(a,Pb))\bigr).
\]
Coverage and successor separation give $Z(a)\in\{0,1\}$ and
$Z(a)=1\leftrightarrow a=0$.  They similarly give
$U(a)\in\{0,1\}$ and $U(a)=1\leftrightarrow a=1$.
If $b=0$, $E(a,b)=1$ is equivalent to $a=0$.
If $b=Sn$, $E(a,b)=1$ is equivalent to $d(a,n)=1$.
The latter condition implies $a=n+1=Sn=b$ by
\eqref{eq:const-reconstruct}; conversely $a=Sn$ gives $d(Sn,n)=1$.
Thus $E(a,b)$ is Boolean and equals one exactly when $a=b$.
The fixed primitive recursive term
\[
                         \delta_*(a,b)=Z(E(a,b))
\]
is therefore Boolean and equals zero exactly when $a=b$.  This
derivation uses only quantifier-free inductions after they have already
been recovered; it is not part of the induction-free bootstrap.

Consequently the chosen equality characteristic satisfies
\textup{(Split)}.  The rule \textup{(Const)} is iPRA-admissible:
under its recurrence premise, equation induction on
$f(n)=f(0)$ gives the conclusion.  All domains remain the whole
numerical domain and equality is interpreted as actual equality.

For completeness, the interpretation composites require more than
agreement in the standard natural numbers.  In the proposed theory,
recovered equation induction proves uniqueness of every fixed primitive
recursive defining recursion, under arbitrary parameters.  By structural
induction on term syntax, it therefore proves that the round trip of
every $R$ term has the same value as the original.  The same argument in
iPRA treats the round trip of each standard primitive recursive function
symbol.  For the additional primitive $\delta$, both $\delta$ and the
translated $\delta_*$ satisfy \textup{(Split)} in the proposed theory.
Case analysis on the first split shows their equality: if $a=b$, both
are zero; if $a\ne b$, both are one.  Hence
\[
                         \delta(a,b)=\delta_*(a,b)
\]
is provable there.  Equality congruence and structural induction on
formulas now give both formula round trips.  The numerical isomorphisms
have graph $x=y$.

\begin{theorem}[Exact scope of the result]
The presentation in this section is whole-domain bi-interpretable with
full first-order iPRA, with both composite numerical interpretations
provably the identity.  It is therefore equiconsistent with iPRA.
The universal statements are written mathematical proofs, not a new
proof-assistant formalization.
\end{theorem}

\subsection{What this does and does not simplify}

The primitive arithmetic inference now concerns a single numerical term
and an unconditional identity-update recurrence.  It has no discharged
induction hypothesis, no Boolean accumulator in its rule text, and no
arbitrary update template.  The accumulator is an untrusted derived
proof construction.

The separate equality split is essential to the proof given here.
Replacing it by an assertion that an internally defined equality test
is faithful, without proving that assertion in the induction-free base,
would be circular.  Likewise, for a tail-recursive iterator,
$d(a,Sb)=d(Pa,b)$ does not by itself imply
$d(a,Sb)=P(d(a,b))$: the missing commutation is a fusion theorem.

Compared with the repository's scalar basis, this presentation adds a
binary term constructor, a stronger primitive equality split in place
of primitive coverage, and a second local binder for the recursor.
It removes arbitrary induction from the primitive inference interface.
No claim is made that those changes reduce the complete checker-size
metric, surpass the existing compact profiles, or establish an absolute
minimum.  The unanswered next question is whether the equality split
can be derived from a sufficiently small recursion base and
\textup{(Const)} alone.
